\chapter{The Assumption Audit}
\label{ch:assumption-audit}

This chapter derives nothing. Its purpose is to state plainly, in one place,
exactly what Parts I and II have established and exactly what standard
quantum mechanics contains that this book has not yet earned the right to
use. After forty-some pages of successful reconstruction it becomes
psychologically easy to feel that quantum mechanics has been derived. It has
not — not yet — and this chapter exists so that no later chapter can quietly
trade on that feeling. One clarification before the accounting itself: the
ledger below records \emph{logical} dependence, not historical dependence.
An idea may have been known for decades — the Born rule predates every
chapter of this book by a century — while still appearing late here, because
\emph{in this reconstruction} it depends on structure established only in
the preceding chapters. The order of this book is not a claim about the
order in which physics discovered these ideas.

\section{What Has Been Established}

\paragraph{Single-system foundations (Part I).} The probability simplex and
stochastic maps (Chapter~\ref{ch:simplex}); bistochastic matrices as a
polytope generated by permutations, and bistochastic evolution as exactly
majorization-compatible evolution with non-decreasing entropy (Chapter
\ref{ch:birkhoff}); the $n=2$ degeneracy of orthostochastic and unistochastic
matrices, and their strict separation at $n=3$ via an explicit
Hadamard-matrix obstruction (Chapter~\ref{ch:unistochastic}).

\paragraph{Compositional sufficiency (Part I, first axis).} Probability-level
composition discards information present in amplitude-level composition —
but this alone does not select $\C$ over $\R$; a real lift shows the same
obstruction (Chapter~\ref{ch:composition}).

\paragraph{Global realizability (Part I, second axis).} Exact frame closure
around a cycle projects to no corresponding closure relation on observed
transition probabilities (Chapter~\ref{ch:context-graphs}); edgewise real
admissibility does not imply global real admissibility, proved via the
triangle-gluing theorem and its mutually-unbiased corollary
(Chapter~\ref{ch:real-mub-obstruction}), and generalized to arbitrary context
graphs via spanning trees and cycle rank
(\S\ref{sec:spanning-trees}). This is a genuine single-system obstruction to
$\R^2$, not dependent on composite systems.

\paragraph{Extremal reformulation (Part II opening).} The real/complex
minimal witness at $d=2$ is exactly $K_3$, $m_\beta(\R,\C;2)=1$
(Chapter~\ref{ch:minimal-witnesses}) — a restatement of the Part I result in
extremal language, with several genuinely open extremal questions
(higher-$d$ witnesses, quaternionic separation, forbidden subgraphs, density
thresholds) explicitly flagged as unresolved.

\paragraph{System composition (Part II, third axis).} Local tomography,
precisely defined via Hermitian-space dimension counting, holds exactly for
$\C$ and fails in opposite directions for $\R$ and $\Hset$
(Chapter~\ref{ch:local-tomography}) — the classical Araki/Wootters/Hardy
argument, not original to this book. For $\C$, tensor products of local
operator bases literally span the composite Hermitian space; for $\R$, the
missing dimension is identified concretely as $J\otimes J$
(Chapter~\ref{ch:tensor-product}). Positivity is derived operationally, and
the $J\otimes J$ coefficient is shown to support a bounded physical family
$|c|\le1$ (Chapter~\ref{ch:positivity}). Both partial traces annihilate this
coefficient (Chapter~\ref{ch:partial-trace}). The same family is
$\R$-entangled for every $c\ne0$ and $\C$-separable throughout its entire
range, via an explicit decomposition — separability is field-relative, not
intrinsic to a state (Chapter~\ref{ch:entanglement}).

\section{What Remains}
\label{sec:what-remains}

None of the above constitutes a derivation of quantum mechanics. The
following pillars of the standard formalism have been used nowhere in Parts
I--II and are not yet earned:

\begin{itemize}
\item[] \textbf{General complex Hilbert-space state structure.} Everything
proved about $\C$ so far concerns finite dimension, and the sharpest worked
examples are $d=2,3$. The general claim that $n$-outcome quantum systems are
described by rays or density operators on $\C^n$ has not been derived for
general $n$, only motivated dimension-by-dimension.

\item[] \textbf{The Born rule, in full.} This book has been careful from
Chapter~\ref{ch:composition} onward to keep four distinct statements separate
and has proved none of the last three:
\begin{equation}
B_{ij} = |U_{ij}|^2 \quad(\text{proved, Ch.~\ref{ch:unistochastic}}), \qquad
p_i = |\psi_i|^2, \qquad p(i|\psi) = |\langle i|\psi\rangle|^2, \qquad p_i =
\operatorname{tr}(\rho P_i).
\end{equation}
The bistochastic-matrix identity is not the same claim as the state-vector
Born rule, which is not the same claim as the general projective-measurement
rule, which is not the same claim as the density-operator/POVM rule.

\item[] \textbf{Hermitian operators as the general observable class.} Part II
used $\mathrm{Herm}_\C(d)$ as the ambient space for states, motivated by
positivity of probabilities; it has not shown that \emph{observables} in
general (as opposed to states) must be Hermitian, nor derived the spectral
decomposition $A = \sum_i a_i P_i$ and the eigenvalue rule for measurement
outcomes.

\item[] \textbf{Unitary dynamics and the Schrödinger equation.} Nothing in
this book has yet introduced continuous time evolution. The one-parameter
group law $U(t+s)=U(t)U(s)$ and the generator theorem $U(t) = e^{-itH/\hbar}
\Rightarrow i\hbar\,\partial_t|\psi\rangle = H|\psi\rangle$ remain to be
derived, not assumed.

\item[] \textbf{General measurement and POVM structure.} Only rank-one
projective tests have been used, to motivate positivity
(Chapter~\ref{ch:positivity}). The general POVM formalism
$p_i=\operatorname{tr}(\rho E_i)$, $\sum_i E_i = I$, $E_i \ge 0$, has not been
derived or even stated as a target.
\end{itemize}

\section{Roadmap for Part III}

Part III attacks these in dependency order. The Born rule is addressed first
(Chapter~\ref{ch:born-rule}), since the composite-system machinery of Part II
already distinguishes exactly the four formulas above and Part III's task is
to reconcile them rather than introduce new ones. Observables and spectral
structure follow, motivated by the same admissibility logic used for states
in Chapter~\ref{ch:positivity}. Continuous reversible composition is then
used to derive the unitary group and Schrödinger's equation as a generator
theorem, in the spirit of Chapter~\ref{ch:context-graphs}'s treatment of
discrete composition. POVMs are introduced last, as the generalization of
projective measurement forced by considering measurements performed on a
subsystem of a larger composite system — connecting directly back to the
partial trace of Chapter~\ref{ch:partial-trace}.

\begin{ledgerentry}[Opening of Part III]
\textbf{Established:} nothing new — this chapter is a compilation of Parts
I--II, cross-referenced above, with no additional claims.\\
\textbf{Not established:} every item in \S\ref{sec:what-remains}. Part III
exists to close these gaps one at a time, each under the same
reconstruction-vs-recovery discipline used throughout: what is derived from
what has already been proved, versus what is imported to check the result
against the standard formalism of Sakurai--Napolitano~\cite{sakurainapolitano}
and Gross~\cite{gross}.
\end{ledgerentry}
